%% 8장 부동소수점 수
\chapter{부동소수점 수}
\label{ch:floating-point}

\abstract{컴퓨터가 실수를 다루는 방법은 거의 예외 없이 IEEE~754 표준의 부동소수점 수이다. 부동소수점 수는 실수의 유한한 부분집합이고, 계산의 결과가 그 집합에 없으면 가장 가까운 원소로 반올림된다. 이 장에서는 부동소수점 수의 집합이 어떤 모양인지, 곧 수들 사이의 간격이 어떻게 변하고 가장 크고 작은 수가 얼마인지를 정리한 뒤, binary64 형식의 비트 배치와 반올림의 규칙을 엄밀하게 정의한다. 이어서 반올림 오차의 기본 정리, 단위 ULP와 그것으로 재는 오차, 뺄셈이 정확해지는 Sterbenz의 조건, 서브노멀 수가 보장하는 점진적 언더플로를 증명하고, 부호 있는 영과 무한대와 NaN이 계산에서 어떻게 행동하는지를 살핀다. 마지막으로 정수와 부동소수점 수 사이의 변환과, 컴파일러가 부동소수점 식을 함부로 바꾸지 못하게 하는 이유를 본다.}

계산기에 $0.1+0.2$를 넣으면 $0.3$이 나오지만, Python에서 \code{0.1 + 0.2 == 0.3}은 \code{False}이고 \code{0.1 + 0.2}는 \code{0.30000000000000004}로 찍힌다. 버그처럼 보이지만 이것은 IEEE~754가 정한 그대로의 동작이며, 이 장을 다 읽고 나면 이 결과를 손으로 예측할 수 있게 된다. $0.1$은 이진법으로 유한하게 적히지 않으므로 \code{0.1}이라는 리터럴은 $0.1$에 가장 가까운 binary64 수를 가리키고, 그 수는 정확히
\[
0.1000000000000000055511151231257827021181583404541015625
\]
이다. \code{0.2}와 \code{0.3}도 마찬가지로 가장 가까운 수로 바뀌고, 두 근삿값을 더한 정확한 합이 다시 가장 가까운 수로 반올림되는데, 그 수가 하필 \code{0.3}이 가리키는 수의 바로 다음 수이다. 모든 단계가 ``가장 가까운 수''라는 하나의 규칙을 따랐을 뿐인데 결과가 어긋난 것이다. 수치 계산의 정확도를 따지는 모든 논의는 이 규칙을 정확히 이해하는 데서 시작한다.

\section{부동소수점 수의 집합}
\label{sec:fp-set}

\subsection{정의}

과학적 표기법은 수를 $6.022\times10^{23}$처럼 유효숫자와 지수로 나누어 적는다. 유효숫자의 자릿수를 고정하면 아주 큰 수와 아주 작은 수를 같은 상대 정밀도로 적을 수 있는데, 부동소수점 수는 이 생각을 기수 $2$로 옮긴 것이다.

\begin{definition}[부동소수점 수]
\label{def:fp-set}
정수 $\beta\ge2$(기수), $p\ge1$(정밀도), $e_{\min}<e_{\max}$에 대해
\[
x=\pm\,m\cdot\beta^{e-p+1},\qquad m\in\mathbb Z,\ 0\le m<\beta^p,\ e_{\min}\le e\le e_{\max}
\]
꼴로 쓸 수 있는 실수 $x$들의 집합을 $\mathbb F$라 한다. $m$을 \term{유효숫자}{significand}, $e$를 \term{지수}{exponent}라 한다. $m\ge\beta^{p-1}$로 쓸 수 있는 $x$를 \term{정규 수}{normal number}, $0$이 아니면서 그렇게 쓸 수 없는 $x$를 \term{서브노멀 수}{subnormal number}라 한다.
\end{definition}

유효숫자 $m$을 $\beta$진법 $p$자리로 적고 맨 앞자리 뒤에 소수점을 찍으면 $x=\pm d_0.d_1d_2\cdots d_{p-1}\times\beta^e$가 되며, 정규 수는 맨 앞자리 $d_0$이 $0$이 아닌 수이다. 서브노멀 수는 지수가 이미 가장 작은 $e_{\min}$이라 앞자리의 $0$을 지수로 옮겨 없앨 수 없는 수이다. 이 책에서 다루는 것은 거의 언제나 $\beta=2$, $p=53$, $e_{\min}=-1022$, $e_{\max}=1023$인 \textbf{binary64}이고, C와 C++의 \code{double}, Python의 \code{float}, numpy와 stardust의 \code{float64}가 모두 이것이다\cite{ieee754-2019}.

\subsection{간격의 구조}

$\mathbb F$의 수들은 고르게 퍼져 있지 않다. 이진법에서 지수가 $e$인 정규 수는 $2^e\le|x|<2^{e+1}$에 놓이는데, 이 구간을 \term{이진 구간}{binade}이라 부르자. 한 이진 구간 안에서는 유효숫자 $m$이 $2^{p-1}$부터 $2^p-1$까지 $1$씩 커지므로 이웃한 두 수의 간격은 일정하게 $2^{e-p+1}$이고, 이진 구간을 하나 올라갈 때마다 간격이 두 배가 된다. 그래서 수들은 $0$ 근처에 촘촘하고 멀어질수록 성기다.

간격을 $x$의 크기로 나눈 상대 간격은 이진 구간의 아래 끝 $2^e$에서 $2^{1-p}$이고 위 끝 근처에서 그 절반 $2^{-p}$에 가깝다. 상대 정밀도가 한 이진 구간 안에서 두 배까지 출렁이는 이 현상을 \term{흔들림}{wobble}이라 하는데, 기수가 클수록 흔들림이 커지므로(기수 $\beta$에서는 $\beta$배) 기수 $2$가 상대 정밀도를 가장 고르게 유지한다. 이것이 이진 부동소수점을 쓰는 이유 가운데 하나이다\cite{goldberg1991}.

\begin{proposition}
\label{prop:fp-integers}
binary64에서 $|n|\le2^{53}$인 모든 정수 $n$은 정확히 나타난다. $2^{53}+1$은 나타나지 않는다.
\end{proposition}

\begin{proof}
$|n|<2^{53}$이면 $m=|n|$, $e=p-1=52$로 두면 $n=\pm m\cdot2^0$이다. $n=\pm2^{53}$은 $m=2^{52}$, $e=53$으로 쓴다. 한편 $2^{53}\le x<2^{54}$인 이진 구간의 간격은 $2^{53-52}=2$이므로 이 구간의 수는 모두 짝수이고, 홀수 $2^{53}+1$은 그 가운데 없다.
\end{proof}

$2^{53}=9007199254740992$이다. 이 경계는 정수와 부동소수점 수를 섞어 쓸 때 자주 만나는데, 64비트 정수의 범위는 $2^{63}$까지이므로 큰 정수를 \code{float64}로 바꾸면 끝자리가 사라질 수 있다.

\subsection{binary64의 비트 배치}

binary64 수 하나는 $64$비트에 다음처럼 담긴다. 맨 위 $1$비트는 부호 $s$, 그다음 $11$비트는 \term{치우친 지수}{biased exponent} $E$, 나머지 $52$비트는 \term{가수 필드}{fraction field} $f$이다. 해석은 $E$에 따라 셋으로 나뉜다.
\begin{itemize}
\item $1\le E\le2046$이면 정규 수 $(-1)^s\,(1+f/2^{52})\cdot2^{E-1023}$이다. 정규 수의 맨 앞자리는 이진법에서 언제나 $1$이므로 그 비트는 저장하지 않으며, 이 \term{숨은 비트}{hidden bit} 덕분에 $52$비트로 $53$자리의 정밀도를 얻는다.
\item $E=0$이면 $(-1)^s\,(f/2^{52})\cdot2^{-1022}$이다. $f=0$이면 $\pm0$이고, 그 밖에는 서브노멀 수이다.
\item $E=2047$이면 $f=0$일 때 $\pm\infty$, $f\ne0$일 때 NaN이다.
\end{itemize}
지수에서 $1023$을 빼는 것은 지수를 부호 없는 정수로 적기 위한 장치이다. 그 결과로 다음의 중요한 성질이 생기는데, 1장에서 부동소수점 수를 정수 키로 정렬할 때 가정으로 받아들였던 것이 바로 이것이다.

\begin{proposition}
\label{prop:fp-bits-monotone}
NaN이 아닌 binary64 수 $x$의 하위 $63$비트, 곧 $E$와 $f$를 이어 붙인 비트열을 부호 없는 정수로 읽은 값을 $\mu(x)=2^{52}E+f$라 하자. NaN이 아닌 두 수 $x$, $y$에 대해 $|x|<|y|$와 $\mu(x)<\mu(y)$는 동치이다. 또한 $\mu$가 $1$ 커지는 것은 $|x|$가 바로 다음 binary64 수로 가는 것과 같다.
\end{proposition}

\begin{proof}
$\mu$를 $0,1,2,\dots$로 늘려 가며 나타내는 수를 보자. $E=0$인 동안에는 $|x|=f\cdot2^{-1074}$이 $f$에 비례해 $2^{-1074}$씩 커진다. $f=2^{52}-1$에서 $f=0$, $E=1$로 넘어가면 $|x|$는 $(2^{52}-1)2^{-1074}$에서 $2^{-1022}=2^{52}\cdot2^{-1074}$로 역시 $2^{-1074}$만큼 커진다. $E\ge1$인 이진 구간 안에서는 $|x|=(2^{52}+f)\cdot2^{E-1075}$가 $f$에 대해 일정한 간격 $2^{E-1075}$로 커지고, $f=2^{52}-1$에서 다음 $E$로 넘어갈 때도 $(2^{53}-1)2^{E-1075}$에서 $2^{53}\cdot2^{E-1075}=2^{52}\cdot2^{E+1-1075}$로 같은 간격만큼 커진다. 마지막 유한 수 다음의 $E=2047$, $f=0$은 $+\infty$이다. 따라서 $\mu\mapsto|x|$는 강한 증가함수이고 연속한 $\mu$는 이웃한 수에 대응한다.
\end{proof}

이 명제에서 부호를 붙여 수를 정수 키로 보내는 1장의 사상이 나왔다. 또 이웃한 수로 가는 연산이 비트 패턴에 $\pm1$을 하는 것으로 줄어든다. 양수를 $+\infty$ 쪽으로 옮기려면 $1$을 더하고, 음수는 같은 방향으로 가려면 크기가 줄어야 하므로 $1$을 뺀다.

\begin{instardust}[\code{csrc/array/kernels/scalar\_ops.hpp}]
\code{nextafter(a, b)}는 \cref{prop:fp-bits-monotone}을 그대로 쓴다. NaN이나 $a=b$ 같은 경우를 먼저 처리하고, $a=0$이면 $b$의 부호를 가진 가장 작은 서브노멀 수를 낸 다음, 나머지 경우에는 $a$의 비트 패턴을 $64$비트 정수로 읽어 $0$에서 멀어지는 방향이면 $1$을 더하고 가까워지는 방향이면 $1$을 뺀다. 주석에 따르면 이 방법은 원소마다 C 라이브러리의 함수를 부르는 numpy보다 약 $8$배 빠르다.
\end{instardust}

$E$와 $f$의 범위를 대입하면 binary64의 주요 상수가 나온다. \cref{tab:binary64}에 적은 값들은 모두 정확한 유리수 산술로 확인한 것이다.

\begin{table}[t]
\begin{sdtable}{colspec={lll}, row{1}={font=\small\bfseries}}
\toprule
이름 & 정확한 값 & 십진 근삿값 \\
\midrule
가장 큰 유한 수 & $(2-2^{-52})\cdot2^{1023}$ & $1.7976931348623157\times10^{308}$ \\
가장 작은 양의 정규 수 & $2^{-1022}$ & $2.2250738585072014\times10^{-308}$ \\
가장 작은 양의 서브노멀 수 & $2^{-1074}$ & $4.9406564584124654\times10^{-324}$ \\
기계 엡실론 $\varepsilon$ & $2^{-52}$ & $2.220446049250313\times10^{-16}$ \\
단위 반올림 $u$ & $2^{-53}$ & $1.1102230246251565\times10^{-16}$ \\
\bottomrule
\end{sdtable}
\caption{binary64의 주요 상수. 기계 엡실론은 $1$과 그다음 수의 간격이고, 단위 반올림은 그 절반이다.}
\label{tab:binary64}
\end{table}

\section{반올림}
\label{sec:rounding}

\subsection{반올림 함수}

실수 $x$를 $\mathbb F$의 원소로 바꾸는 함수를 반올림이라 한다. IEEE~754의 기본값은 가장 가까운 원소로 보내되, 두 원소의 한가운데에 있으면 유효숫자가 짝수인 쪽, 곧 마지막 비트가 $0$인 쪽으로 보내는 \term{짝수로의 반올림}{round to nearest, ties to even}이다. 이 밖에 $0$ 쪽, $+\infty$ 쪽, $-\infty$ 쪽으로 자르는 세 가지 \term{방향 반올림}{directed rounding}이 있는데, 방향 반올림은 구간 산술처럼 참값을 위아래로 확실히 묶어야 하는 곳에서 쓰인다. 별다른 말이 없으면 이 책에서 반올림은 짝수로의 반올림이고 이것을 $\fl(x)$로 적는다.

반올림을 정의할 때는 $\mathbb F$의 범위 밖에서 무슨 일이 일어나는지도 정해야 한다. 표준은 지수의 상한이 없는 가상의 집합에서 반올림한 결과의 크기가 가장 큰 유한 수를 넘으면 $\pm\infty$를 결과로 하는데, 짝수로의 반올림에서는 이것이 $|x|\ge(2-2^{-53})\cdot2^{1023}$, 곧 가장 큰 유한 수와 그다음의 가상의 수 $2^{1024}$의 한가운데 이상일 때이다. 이것을 \term{넘침}{overflow}이라 한다.

짝수 쪽으로 가르는 규칙이 하필 필요한 이유는 한가운데 값이 생각보다 자주 나오기 때문이다. 두 수를 더하거나 정수를 바꿀 때 버리는 부분이 정확히 마지막 자리의 절반인 경우는 드물지 않은데, 그때마다 늘 위로 올리면 오차가 한 방향으로 쌓인다. 짝수로 가르면 한가운데 값의 절반 정도는 위로, 절반 정도는 아래로 가므로 이런 치우침이 사라진다. Knuth는 이 규칙에서 $(u\oplus v)\ominus v$를 거듭해도 값이 첫 단계 이후로는 더 바뀌지 않는다는 것을 보였는데, 늘 위로 올리는 규칙에서는 같은 계산이 한 방향으로 끝없이 흘러갈 수 있다\cite{knuth1997taocp2}. 정수 변환에서 짝수 쪽으로 가르는 예를 보면, $2^{53}+1$은 $2^{53}$과 $2^{53}+2$의 한가운데이고 유효숫자가 짝수인 쪽은 $2^{53}$이므로 $\fl(2^{53}+1)=2^{53}$이다. 반면 $2^{53}+3$은 $2^{53}+2$와 $2^{53}+4$의 한가운데이고 $2^{53}+4$ 쪽이 짝수이므로 위로 올라간다.

\subsection{올바르게 반올림된 연산}

IEEE~754의 핵심 약속은 덧셈, 뺄셈, 곱셈, 나눗셈, 제곱근, 그리고 곱한 뒤 더하는 fused multiply-add가 \term{올바르게 반올림된다}{correctly rounded}는 것이다. 곧 두 수 $a,b\in\mathbb F$에 대해 컴퓨터가 내는 $a\oplus b$는 정확한 실수 $a+b$를 반올림한 $\fl(a+b)$와 같아야 하고, 다른 연산도 마찬가지이다\cite{ieee754-2019}. 이 약속 덕분에 같은 연산이 어느 컴퓨터에서나 같은 비트를 내고, 오차의 분석이 연산마다 반올림 한 번으로 줄어든다. 반면 $\exp$나 $\sin$ 같은 초월함수에는 이 약속이 없다. 올바른 반올림을 하려면 결과가 두 수의 한가운데에 아주 가까운 경우를 가려내야 하는데 그 비용을 정하기 어렵기 때문이며, 그래서 라이브러리마다 초월함수의 마지막 비트가 다를 수 있다. 이 문제는 11장에서 다룬다.

반올림의 오차는 다음 정리로 요약된다. 정리를 깔끔하게 적으려고 범위를 정규 수의 영역으로 한정한다.

\begin{theorem}[반올림 오차의 기본 정리]
\label{thm:standard-model}
$x$가 실수이고 $2^{-1022}\le|x|\le\max\mathbb F$이면
\[
\fl(x)=x(1+\delta),\qquad|\delta|\le u=2^{-53}
\]
이다. 또한 $\fl(x)=x/(1+\delta')$, $|\delta'|\le u$로도 쓸 수 있다.
\end{theorem}

\begin{proof}
$2^e\le|x|<2^{e+1}$이라 하자($e\ge-1022$). 이 이진 구간의 이웃한 수들의 간격은 $2^{e-52}$이므로 가장 가까운 수까지의 거리는 그 절반 이하이고, $|\fl(x)-x|\le2^{e-53}$이다. 위 끝에서 $\fl(x)=2^{e+1}$로 반올림되는 경우에도 이 부등식은 그대로이다. 따라서 $|\delta|=|\fl(x)-x|/|x|\le2^{e-53}/2^e=2^{-53}$이다. 둘째 꼴은 $\delta'=(x-\fl(x))/\fl(x)$로 두면 되는데, $|\fl(x)|\ge2^e$이므로 같은 계산으로 $|\delta'|\le2^{-53}$이다.
\end{proof}

연산이 올바르게 반올림되므로 이 정리는 곧바로 연산의 오차가 된다. 넘침과 언더플로가 없으면 $a\oplus b=(a+b)(1+\delta)$, $|\delta|\le u$이고 다른 연산도 같은 꼴이다. 이것을 \term{표준 모형}{standard model}이라 부르며, 9장의 반올림 오차 해석은 모두 이 한 줄 위에 세워진다. 정리의 등호 조건도 짚어 두자. 상대오차는 $x$가 이진 구간의 아래 끝 바로 위에 있으면서 이웃한 두 수의 한가운데에 있을 때 $u$에 가장 가깝고, 위 끝 근처에서는 절반 정도로 줄어든다. 흔들림이 오차의 한계에도 나타나는 것이다.

서막의 예로 돌아가자. $0.1$의 근삿값의 상대오차는 약 $5.55\times10^{-17}$로 $u/2$ 정도이다. \code{0.1}과 \code{0.2}의 근삿값을 정확히 더하면 $0.3000000000000000166\cdots$인데, $[0.25,0.5)$의 간격 $2^{-54}$로 잰 이 값의 이웃은 \code{0.3}의 근삿값 $0.2999999999999999888\cdots$과 그다음 수 $0.3000000000000000444\cdots$이다. 정확한 합은 두 수의 정확히 한가운데에 있고, 짝수로 가르는 규칙이 뒤의 수를 고른다. 비트 패턴으로는 \code{0x3FD3333333333333}과 \code{0x3FD3333333333334}의 차이이다.

\subsection{ULP}

반올림 오차를 상대오차 대신 ``마지막 자리 몇 개''로 재는 것이 편할 때가 많다.

\begin{definition}[ULP]
\label{def:ulp}
$2^e\le|x|<2^{e+1}$인 실수 $x$에 대해 $\operatorname{ulp}(x)=2^{\max(e,-1022)-52}$로 정의하고, $\operatorname{ulp}(0)=2^{-1074}$로 둔다.
\end{definition}

$\operatorname{ulp}(x)$는 $x$가 속한 이진 구간에서 이웃한 binary64 수의 간격이며, \cref{thm:standard-model}의 증명은 곧 $|\fl(x)-x|\le\operatorname{ulp}(x)/2$를 보인 것이다. ULP의 정의는 문헌마다 조금씩 다른데, 차이는 $x$가 $2$의 거듭제곱에 걸릴 때 어느 쪽 간격을 쓰느냐에 있으며 여기의 정의는 Goldberg의 것에 해당한다\cite{goldberg1991,muller2018handbook}. 근삿값 $\hat y$와 참값 $y$의 오차를 $|\hat y-y|/\operatorname{ulp}(y)$로 재면 상대오차보다 흔들림의 영향이 적고 정수 근처의 값으로 나오므로 해석이 쉽다. ``$1$ ULP 이내''라는 말은 대략 상대오차 $2^{-52}$에서 $2^{-53}$ 이내라는 뜻이다.

두 binary64 수 사이의 거리를 ULP로 잴 때는 더 단순한 방법이 있다. \cref{prop:fp-bits-monotone}과 1장의 정수 키 $\kappa$에서, 두 수 $a$와 $b$ 사이에 있는 binary64 수의 개수는 $|\kappa(a)-\kappa(b)|$이다. $0$을 사이에 둔 두 수에 대해서는 양쪽의 아주 작은 서브노멀 수들까지 모두 세어지는 점이 상대오차와 다르다.

\begin{instardust}[\code{bench/accuracy/common.py}]
stardust의 정확도 측정은 \code{ulp\_distance}로 두 값 사이의 binary64 수의 개수를 센다. 두 값을 순서를 보존하는 정수로 바꾼 뒤 큰 쪽에서 작은 쪽을 빼는데, $-\infty$에서 $+\infty$까지의 거리는 부호 있는 $64$비트 정수에 들어가지 않으므로 부호 없는 정수의 뺄셈을 쓴다. 2장에서 보았듯이 이 뺄셈은 $2^{64}$의 나머지로 감기지만 참값이 $2^{64}$보다 작으므로 결과는 정확하다. $+0$과 $-0$의 거리는 $0$이고 NaN 둘의 거리도 $0$으로 둔다. 각 영역의 정확도 약속은 ``Arb로 구한 참값을 binary64로 반올림한 값에서 $k$ ULP 이내''의 꼴로 적힌다.
\end{instardust}

\subsection{정확한 연산}

반올림이 늘 오차를 만드는 것은 아니며, 결과가 이미 $\mathbb F$에 있으면 연산은 정확하다. $2$의 거듭제곱을 곱하거나 나누는 것은 넘침과 언더플로가 없으면 지수만 바꾸므로 정확하고, \cref{prop:fp-integers}의 범위 안에서 정수끼리의 덧셈과 곱셈도 결과가 $2^{53}$ 이하이면 정확하다. 이보다 덜 자명한 정확성의 원천은 가까운 두 수의 뺄셈이다.

\begin{lemma}[Sterbenz]
\label{lem:sterbenz}
$x,y\in\mathbb F$가 $y/2\le x\le2y$를 만족하면 $x-y\in\mathbb F$이다. 따라서 $x\ominus y=x-y$는 정확하다.
\end{lemma}

\begin{proof}
$x$와 $y$는 같은 부호이므로 둘 다 양수라 해도 된다. 대칭성에서 $y\le x\le2y$라 하자. $y$의 지수를 $e$라 하면 $x$와 $y$는 모두 $2^{\max(e,-1022)-52}$의 정수배이다. $y$는 그 정의에서, $x$는 $x\ge y$이므로 지수가 $e$ 이상이어서 간격이 $2^{\max(e,-1022)-52}$의 배수이기 때문이다. 따라서 $x-y$도 이 수의 정수배이다. 한편 $0\le x-y\le y<2^{e+1}$이므로 $x-y$는 이 간격의 $2^{53}$배 미만이고, 유효숫자 $53$비트로 적힌다. 그러므로 $x-y\in\mathbb F$이다.
\end{proof}

정리의 조건은 두 수가 두 배 이내로 가깝다는 것뿐이다. 가까운 수의 뺄셈은 상대오차를 크게 키우는 \textbf{소거}의 원인으로 악명이 높은데, Sterbenz의 보조정리는 소거가 뺄셈 자체에서는 오차를 만들지 않는다는 것을 말해 준다. 소거가 위험한 것은 뺄셈이 \textbf{이미 오차를 품은} 두 값의 앞자리를 지워서 그 오차를 드러내기 때문이며, 9장에서 이 구별을 정식으로 다룬다. 정확한 뺄셈은 그 자체로 쓸모가 있어서, 오차 없는 변환(10장)과 초월함수의 인자 축소(11장)에서 거듭 쓰인다.

\begin{instardust}[\code{docs/design/array.md} §5.4]
배열의 \code{expm1} 커널은 $e^x=s\cdot(1+t)$로 나눈 뒤 $e^x-1$을 \code{fma(s, t, s - 1)}로 계산한다. $s\in[1/2,2]$이므로 \cref{lem:sterbenz}에서 $s-1$은 정확하고, 반올림은 마지막 fused multiply-add 한 번뿐이다. 설계 문서는 이 근거를 ``Sterbenz로 정확하고''라고 적어 둔다.
\end{instardust}

\section{서브노멀 수와 점진적 언더플로}
\label{sec:subnormal}

정규 수만 있다면 $0$과 가장 작은 정규 수 $2^{-1022}$ 사이에 큰 빈틈이 생긴다. 그 빈틈의 너비는 $2^{-1022}$이지만 바로 위의 간격은 $2^{-1074}$이므로, 빈틈이 이웃한 간격보다 $2^{52}$배나 넓다. 계산 결과가 이 빈틈에 떨어지면 $0$으로 반올림하는 수밖에 없고, 이것을 \term{갑작스런 언더플로}{abrupt underflow} 또는 영으로의 플러시라 한다. 서브노멀 수는 이 빈틈을 정규 수의 마지막 간격 $2^{-1074}$로 고르게 채워서, 값이 작아질 때 정밀도가 한꺼번에 사라지지 않고 비트 하나씩 줄어들게 한다. 이것이 \term{점진적 언더플로}{gradual underflow}이다.

\begin{proposition}
\label{prop:gradual-underflow}
서브노멀 수가 있는 binary64에서 $x,y\in\mathbb F$가 유한하고 $x\ne y$이면 $x\ominus y\ne0$이다. 서브노멀 수를 $0$으로 플러시하면 이 성질이 깨진다.
\end{proposition}

\begin{proof}
모든 유한한 binary64 수는 $2^{-1074}$의 정수배이므로 $x-y$도 $0$이 아닌 $2^{-1074}$의 정수배이다. $|x-y|<2^{-1022}$이면 $x-y=k\cdot2^{-1074}$, $0<|k|<2^{52}$이므로 $x-y$는 서브노멀 수로 정확히 나타나고 $x\ominus y=x-y\ne0$이다. $|x-y|\ge2^{-1022}$이면 반올림 결과의 크기도 $2^{-1022}$ 이상이다. 플러시하는 경우의 반례로 $x=1.5\cdot2^{-1022}$, $y=2^{-1022}$를 들 수 있는데, 두 수는 모두 정규 수이지만 차 $2^{-1023}$은 서브노멀 수이므로 $0$이 된다.
\end{proof}

이 성질이 없으면 \code{if (x != y) z = 1 / (x - y);} 같은 코드가 $0$으로 나눌 수 있다. 점진적 언더플로가 있으면 표준 모형도 언더플로 영역까지 넓어진다. 반올림 결과가 서브노멀 영역에 떨어지면 상대오차는 더 이상 $u$로 묶이지 않지만 절대오차는 간격의 절반 $2^{-1075}$로 묶이므로, 모든 유한한 결과에 대해
\[
\fl(x)=x(1+\delta)+\eta,\qquad|\delta|\le u,\quad|\eta|\le2^{-1075},\quad\delta\eta=0
\]
이 성립한다. 곱셈과 나눗셈에는 이 꼴이 필요하지만, 덧셈과 뺄셈은 \cref{prop:gradual-underflow}의 증명에서 보았듯 결과가 서브노멀 영역에 떨어지면 언제나 정확하므로 $\eta=0$으로 둘 수 있다.

\begin{aside}[Kahan과 IEEE~754]
1970년대의 컴퓨터는 회사마다 부동소수점 형식과 반올림 규칙이 달라서, 같은 프로그램이 기계를 옮기면 다른 답을 내는 일이 흔했다. 1977년에 시작된 IEEE의 표준화 위원회에서 William Kahan은 Intel의 8087 보조 프로세서 설계에 참여하며 다듬은 제안을 내놓았는데, 가장 오래 논쟁거리가 된 것이 점진적 언더플로였다. 다른 회사들은 서브노멀 수가 하드웨어를 복잡하게 하고 느리게 만든다며 반대했고, Kahan 쪽은 \cref{prop:gradual-underflow} 같은 성질이 프로그램을 증명할 수 있게 해 준다고 맞섰다\cite{coonen1981underflow}. 결국 1985년의 표준은 Kahan의 안을 거의 그대로 받아들였고, Kahan은 이 공로로 1989년 Turing상을 받았다\cite{severance1998}. 오늘날 거의 모든 범용 프로세서가 이 표준을 따른다.
\end{aside}

\section{특수한 값}
\label{sec:special-values}

\subsection{부호 있는 영}

binary64에는 $+0$과 $-0$이 따로 있다. 두 영은 \code{==}로 같다고 비교되지만 몇몇 연산에서 다르게 행동한다. $1/(+0)=+\infty$이고 $1/(-0)=-\infty$이며, $\sqrt{-0}=-0$이다. 덧셈에서는 $(-0)+(-0)=-0$이지만, 짝수로의 반올림에서 크기가 같고 부호가 다른 두 수의 합은 $+0$이므로 $x-x=+0$이고 $(-0)+(+0)=+0$이다. 따라서 $-0$을 더해도 모든 $x$에서 $x$가 그대로 나오지만 $+0$을 더하면 $-0$이 $+0$으로 바뀐다. 덧셈의 항등원이 $-0$인 셈이다.

부호 있는 영이 있는 이유는 $0$으로 다가가는 방향의 정보를 남기기 위해서이다. 아래로 넘쳐 $0$이 된 음수 $-10^{-400}$은 $-0$이 되므로 그 역수가 $-\infty$로 올바른 부호를 갖는다. 5장에서 본 복소함수의 가지 자름 위의 값도 부호 있는 영으로 구별되는데, 예를 들어 $\sqrt{-1+0i}=i$이고 $\sqrt{-1-0i}=-i$로 두면 자름의 위아래에서 각각의 연속성이 지켜진다\cite{kahan1987branch}. 그래서 부호 있는 영을 다룰 때는 \code{x == 0}으로 부호를 알 수 없다는 점에 주의해야 하며, 부호는 \code{signbit}이나 \code{copysign}으로 읽는다.

\subsection{무한대}

$\pm\infty$는 넘침의 결과이자 $1/0$처럼 극한이 정해진 연산의 결과이다. 무한대와의 산술은 극한의 규칙을 따라 $\infty+1=\infty$, $2\cdot\infty=\infty$, $1/\infty=0$이고, 비교에서는 모든 유한한 수보다 크거나 작다. 무한대 덕분에 넘침이 일어나도 계산이 멈추지 않고 이어지며, 넘친 중간값이 나중에 $1/\infty=0$처럼 의미 있는 값으로 돌아오기도 한다.

\subsection{NaN}

$0/0$, $\infty-\infty$, $0\cdot\infty$, $\sqrt{-1}$처럼 극한조차 정해지지 않는 연산의 결과는 NaN(Not a Number)이다. 비트로는 $E=2047$이고 $f\ne0$인 패턴이 모두 NaN이므로 NaN은 $2^{53}-2$개나 된다. 가수 필드의 맨 위 비트가 $1$이면 \term{조용한 NaN}{quiet NaN}, $0$이면 \term{신호 NaN}{signaling NaN}이다. 조용한 NaN은 연산에 들어가면 예외 없이 NaN을 내고, 신호 NaN은 연산에 들어가는 순간 무효 연산 예외를 일으킨다. 나머지 비트는 \term{페이로드}{payload}라 하여 NaN이 생긴 원인 같은 정보를 담는 데 쓸 수 있지만 언어와 라이브러리가 이것을 보존한다는 보장은 거의 없다.

NaN의 가장 중요한 성질은 1장에서 보았듯이 어떤 값과도 순서가 없다는 것이다. 그래서 \code{x != x}는 $x$가 NaN일 때만 참이고, NaN이 섞이면 최댓값과 최솟값의 의미도 다시 정해야 한다. IEEE~754-2019는 두 가지를 정한다. \code{maximum}과 \code{minimum}은 인자 하나라도 NaN이면 NaN을 내고 $-0<+0$으로 보며, \code{maximumNumber}와 \code{minimumNumber}는 NaN을 빠진 값처럼 무시하고 다른 인자를 낸다\cite{ieee754-2019}. numpy의 \code{maximum}과 \code{fmax}가 각각 이 두 규칙에 해당한다.

\begin{instardust}[\code{csrc/array/kernels/scalar\_ops.hpp}]
\code{fmax\_nan}과 \code{fmin\_nan}은 IEEE~754-2019의 \code{maximum}과 \code{minimum}이다. 인자 가운데 NaN이 있으면 그것을 내고, 두 인자가 같다고 비교되면(곧 $\pm0$이면) 부호 비트를 보아 최댓값으로 $+0$을, 최솟값으로 $-0$을 고른다. 주석에 적힌 대로 이것은 NEON의 \code{FMAX}와 \code{FMIN} 명령이 하는 일과 같아서, 벡터 커널과 스칼라 정의가 같은 비트를 낸다.
\end{instardust}

\subsection{예외와 기본 결과}

표준은 다섯 가지 \term{예외}{exception}를 정하고, 각 예외에 대해 계산을 멈추지 않고 내놓을 \textbf{기본 결과}를 정한다. 무효 연산은 NaN을, 영으로 나누기는 올바른 부호의 무한대를, 넘침은 무한대를, 언더플로는 서브노멀 수나 $0$을, 부정확은 반올림한 값을 낸다. 각 예외가 일어나면 상태 플래그가 세워지고 프로그램이 나중에 이 플래그를 확인할 수 있지만, 고수준 언어에서 플래그를 들여다보는 일은 드물다. 실제로는 기본 결과 덕분에 계산이 끝까지 이어지고, 결과에 남은 NaN이나 무한대를 보고 문제가 있었음을 알게 된다.

\section{정수와의 변환}
\label{sec:conversion}

정수 $n$을 binary64로 바꾸는 것은 반올림 $\fl(n)$이다. $|n|\le2^{53}$이면 정확하고, 그 위에서는 짝수로의 반올림으로 끝자리가 사라진다. 64비트 정수 배열을 \code{float64}로 바꾸어 계산하는 라이브러리가 흔하므로, 정수의 합이나 평균이 $2^{53}$을 넘는 큰 값을 다룰 때는 이 경계를 의식해야 한다.

반대 방향, binary64를 정수로 바꾸는 일에는 세 가지 문제가 있다. 소수 부분을 어떻게 버릴지, 정수형의 범위를 넘는 값을 어떻게 할지, NaN을 무엇으로 할지이다. C와 C++는 소수 부분을 $0$ 쪽으로 자르도록 정하지만, 결과가 정수형의 범위를 벗어나거나 값이 NaN이면 \textbf{정의되지 않은 동작}으로 둔다\cite{iso14882-2020}. 2장에서 부호 있는 정수의 넘침에서 본 것과 같은 함정으로, 컴파일러는 그런 입력이 없다고 가정하고 코드를 바꿀 수 있다.

\begin{instardust}[\code{docs/design/array.md} §2.4]
stardust의 \code{astype}은 \code{float64}를 \code{int64}로 바꿀 때 $0$ 쪽으로 자르고, 범위 밖의 값은 가장 크거나 작은 정수로 포화시키며, NaN은 $0$으로 보낸다. ARM의 \code{FCVTZS} 명령 하나가 정확히 이 일을 하므로 C++의 평범한 캐스트 대신 이 명령을 직접 부른다. 세 경우를 분기로 적으면 정의된 동작은 얻지만 루프가 벡터화되지 않는다는 점도 이 선택의 이유이다. 한편 정수 배열의 \code{mean}은 원소를 \code{float64}로 바꾸어 더하지 않고 정확한 정수합을 구한 뒤 마지막에 한 번만 반올림하므로, 합이 $2^{53}$을 넘어도 결과가 반올림 한 번의 오차 안에 있다.
\end{instardust}

반올림 함수 \code{round}의 동작도 언어마다 다르다는 점을 덧붙인다. Python의 \code{round}와 numpy의 \code{round}는 한가운데 값을 짝수로 보내서 \code{round(2.5)}가 $2$이지만, C의 \code{round}는 $0$에서 먼 쪽으로 보내서 $3$을 낸다. 같은 이름의 함수가 서로 다른 반올림 규칙을 따르는 것이다.

\section{컴파일러와 부동소수점 식}
\label{sec:compiler-fp}

실수의 덧셈은 결합법칙을 만족하지만 부동소수점 덧셈은 그렇지 않다. $(2^{53}+1)+1$을 왼쪽부터 계산하면 첫 덧셈의 결과가 짝수로 반올림되어 $2^{53}$이 되고 다시 $1$을 더해도 $2^{53}$이지만, $2^{53}+(1+1)$은 정확히 $2^{53}+2$이다. 그래서 컴파일러는 부동소수점 식의 계산 순서를 바꿀 수 없고, 프로그래머가 적은 순서대로 연산을 해야 한다.

많은 컴파일러에는 이 제약을 푸는 선택지가 있다. 이른바 \code{-ffast-math}는 결합법칙과 분배법칙이 성립한다고 가정하고, NaN과 무한대가 나오지 않는다고 가정하며, 부호 있는 영을 구별하지 않는다. 덕분에 합을 여러 갈래로 나누어 벡터화하는 등 더 빠른 코드를 만들 수 있지만, 그 가정이 깨지는 입력에서는 결과를 예측할 수 없다. 예를 들어 NaN이 나오지 않는다고 가정한 컴파일러는 \code{x != x} 검사를 거짓으로 바꾸어 없애 버릴 수 있고, 10장의 보상 합처럼 반올림 오차를 정확히 계산해 내는 알고리즘은 대수적으로 $0$인 식을 지워 버리는 최적화 앞에서 무력해진다.

이보다 미묘한 것이 \term{축약}{contraction}이다. 많은 프로세서에는 $a\cdot b+c$를 반올림 한 번으로 계산하는 fused multiply-add(FMA) 명령이 있고, 이것은 곱을 먼저 반올림하고 다시 더하는 것보다 정확하다. 그래서 C 표준은 컴파일러가 \code{a * b + c}를 FMA로 축약하는 것을 허락한다. 하지만 축약은 결과의 비트를 바꾸고, 같은 식이 인라인된 위치나 벡터화 여부에 따라 축약되기도 하고 되지 않기도 하므로, 같은 수식이 경로마다 다른 비트를 낼 수 있다.

\begin{instardust}[\code{CMakeLists.txt}]
stardust는 \code{-ffast-math}를 쓰지 않는다. 빌드 설정의 주석은 ``NaN과 inf는 그 의미를 지켜야 한다''고 이유를 밝힌다. 또 \code{-ffp-contract=off}로 컴파일러의 FMA 축약을 끄는데, 그대로 두면 같은 공식이 다른 경로에서 다른 비트를 낼 수 있기 때문이다. 예컨대 배열 $x$에 함수를 적용한 결과의 $i$번째 원소와 $x[i]$ 하나에 적용한 결과가 같아야 하고, NEON 벡터 커널이 그 스칼라 정의와 같아야 한다. FMA가 필요한 곳에서는 코드가 \code{std::fma}나 \code{vfmaq\_f64}를 직접 적는다. 같은 이유로 컴파일러가 같은 인자의 \code{sin}과 \code{cos}을 한 번의 호출로 합치는 것도 막고, 합치고 싶은 곳에서는 코드가 직접 합친다.
\end{instardust}

\section{맺음}

부동소수점 수는 실수를 흉내 내는 불완전한 장치가 아니라 정확한 규칙을 가진 하나의 수 체계이다. 수들의 집합과 간격이 정해져 있고, 기본 연산은 정확한 결과를 그 집합으로 반올림한 것이며, 그 반올림의 오차는 \cref{thm:standard-model}로 묶인다. 특수한 값들도 각자의 규칙을 따르므로, 계산의 결과는 원칙적으로 언제나 예측할 수 있다. 다음 장에서는 이 규칙을 출발점으로 삼아 연산이 여러 번 이어질 때 오차가 어떻게 쌓이고 퍼지는지를 해석하고, 10장에서는 반올림 오차 자체를 정확히 계산해 내는 방법으로 정확도를 되찾는다.

부동소수점 산술 전반에 대해서는 Goldberg의 개관\cite{goldberg1991}이 고전적인 입문이고, \cite{overton2001}은 짧고 친절한 교과서이며, \cite{muller2018handbook}은 이 분야의 표준적인 참고서이다. Sterbenz의 보조정리는 \cite{sterbenz1974}에서 처음 정리되었다.

\chapterreferences
